% a01_H0rel part 6x1, continued: two-slot bounds of the box propagator at a constant field (input from a01_p6x1.tex after prop:a01:P12)
\begin{lemma}[two-slot bounds of the box propagator at a constant field]\label{lem:a01:RBK}
Let $i\ge1$, $\eta\defeq L^{-i}$ the fine spacing in scale-$i$ units, $\square$ a box of large blocks (a product of intervals, each a union of large blocks) that does not wrap around the torus, with Neumann conditions on its faces, $B_0$ a constant field, and $G_0\defeq G_i(\square,B_0)$ the propagator of step (7) of the proof of \cref{prop:a01:P12}, $G_0^{-1}=D_{B_0}^*D_{B_0}+m^2\ell_i^2+a_iQ_i(B_0)^*Q_i(B_0)$ on $\ell^2(\square\cap\eta\Z^d)$ with the uniform block average $Q_i$. Kernels are taken for the measure $\eta^d$, and $D^a_{B_0,x}$, $a\in\{0,1\}$, denotes the identity or a covariant forward difference on a bond of $\square$ at $x$. \textup{(a)} There are $C,\delta>0$ depending on $(d,L,a)$ and the mass range only, free of $i$, $\eta$, $\square$, $B_0$ and the positions, such that for all $x,x'\in\square$ and $a,b\in\{0,1\}$ ($d\ge3$; here $d=3$)
\[
|D^a_{B_0,x}D^b_{B_0,x'}G_0(x,x')|\ \le\ C\,(\eta+|x-x'|)^{-(d-2+a+b)}\,e^{-\delta|x-x'|},
\qquad\text{so}\qquad |D^a_{B_0,x}D^b_{B_0,x'}G_0(x,x')|\le5C\,e^{-\delta|x-x'|}\ \text{ for }|x-x'|\ge\tfrac35\ (d=3).
\]
\textup{(b)} No depth condition enters: the bound holds at every pair of points of $\square$, the bonds of its faces included.

\textup{(c)} \emph{One slot in Ba{\l}aban's norm.} Let $\Delta,\Delta'$ be two cubes of $\square$ at distance $\ge1$, let $h'$ be a localisation function on $\square$, $0\le h'\le1$, equal to $1$ on the $\eta$-neighbourhood of $\Delta'$ and vanishing at distance $\ge\frac12$ from $\Delta'$, with $|\partial^\eta h'|\le c_h$ and $|\partial^\eta\partial^\eta h'|\le c_h$, and let $x$ lie within $\eta$ of $\Delta$. Then for $a\in\{0,1\}$ the field $\Psi\defeq h'(\cdot)\,D^a_{B_0,x}G_0(x,\cdot)$ satisfies the three clauses of Ba{\l}aban's external-field norm \cite[(1.32)]{Ba3} (sup, first differences, $\alpha_0$-H\"older quotients of the first differences on unit balls) on all of $\square$, its faces included, with the bound $C(1+c_h)^2$, $C$ of the same kind as above.
\end{lemma}
\begin{proof}
(1) \emph{Gauge-out.} With $\lambda_0(x)\defeq B_0\cdot(x-c)$ on the non-wrapping $\square$, $G_i(\square,B_0)=U(\lambda_0)G_i(\square,0)U(\lambda_0)^*$ and $D_{B_0}(U(\lambda_0)\phi)=U(\lambda_0)D\phi$ bond by bond (step (5) of the proof of \cref{prop:a01:P12}; $|U|=1$), so every modulus in the claim equals the one at $B_0=0$. Let $B_0=0$.

(2) \emph{Even images.} Let $\square=\prod_\mu[\alpha_\mu,\beta_\mu)$ and let $T_\square$ be the torus of periods $2(\beta_\mu-\alpha_\mu)$, each a multiple of $2s_b$, so that the large blocks of $\square$ tile $T_\square$ periodically; let $\mathcal W_\square\cong(\Z_2)^d$ be the group generated by the reflections in the face planes $x_\mu=\alpha_\mu-\eta/2$ of $\square$ (planes between sites on block boundaries; on $T_\square$ the reflection in $x_\mu=\beta_\mu-\eta/2$ is the same element). They map sites to sites, bonds to bonds and blocks to blocks, the uniform average commutes with them, and the mass term is constant; every bond of $T_\square$ that crosses a face plane joins a site to its mirror image, where a $\mathcal W_\square$-even field has zero difference, and the remaining bonds of $T_\square$ inside $\square$ are the bonds of $\square$. Hence, exactly as in \cref{lem:a01:PE}(a) and \cref{lem:a01:NL}, the Neumann operator on $\square$ is the even restriction of the periodic one, $G_i(T_\square,0)\Ebar=\Ebar\,G_i(\square,0)$ with $\Ebar$ the even extension, and
\[
G_i(\square,0)(x,x')=\sum_{g\in\mathcal W_\square}G_i(T_\square,0)(x,gx')\qquad(x,x'\in\square),
\]
a difference in $x'$ along a bond of $\square$ acting on the image terms along the reflected bonds. Every image satisfies $\dist_{T_\square}(x,gx')\ge|x-x'|$: coordinatewise, for $x_\mu,x'_\mu\in[\alpha_\mu,\beta_\mu-\eta]$ and $r$ the reflection in $\beta_\mu-\eta/2$, $rx'_\mu-x_\mu=|x_\mu-x'_\mu|+2(\beta_\mu-\tfrac\eta2-\max(x_\mu,x'_\mu))\ge|x_\mu-x'_\mu|$, and the complementary distance around the period is $x_\mu+x'_\mu-2\alpha_\mu+\eta\ge|x_\mu-x'_\mu|$.

(3) \emph{Periodisation.} At zero field $G_i(T_\square,0)(x,y)=\sum_{n\in\Z^d}G_i(\eta\Z^d,0)(x,y+2\ell n)$, $\ell\defeq(\beta_\mu-\alpha_\mu)_\mu$ (the periodic operator is the periodisation of the free-lattice one, as in step (3) of the proof of \cref{prop:a01:P12}), every term having $|x-y-2\ell n|\ge\dist_{T_\square}(x,y)$, and at most $(2+k)^d$ values of $n$ have $|x-y-2\ell n|<k+1$, every period being $\ge1$. So, with the free-lattice bound (4) below at the rate $\delta'$ and $\rho\defeq|x-x'|$,
\[
|D^a_xD^b_{x'}G_i(\square,0)(x,x')|\ \le\ 2^dC'(\eta+\rho)^{-(d-2+a+b)}\sum_{k\ge\lfloor\rho\rfloor}(2+k)^de^{-\delta'k}\ \le\ C(\eta+\rho)^{-(d-2+a+b)}e^{-\delta'\rho/2}.
\]

(4) \emph{The free lattice.} On $\eta\Z^d$ at zero field, $G_i(\eta\Z^d,0)=G_1+\sum_{j=1}^{i-1}G_{(j)}$ with the slices $G_{(j)}\defeq G_{j+1}-G_j=a_jG_jQ^*C^{(j)}QG_j$ (the Schur identity behind \cref{lem:a01:legs}(b) and \cite[(2.17)]{K1}, with $C^{(j)}$ as in \cref{lem:a01:A23}(c)), so that the sum telescopes to $G_i$. In unit-$j$ units ($j$-blocks of side one, spacing $L^{-j}$) let $\ell_j(\cdot,u)\defeq a_jG_jQ^*e_u$ be the leg of the unit datum on the $j$-block $u$; since $G_j$ is real symmetric at zero field, $G_{(j)}(y,y')=a_j^{-1}\sum_{u,v}\ell_j(y,u)C^{(j)}(u,v)\ell_j(y',v)$, each variable sitting in its own leg, so
\[
D^a_yD^b_{y'}G_{(j)}(y,y')=a_j^{-1}\sum_{u,v}(D^a\ell_j)(y,u)\,C^{(j)}(u,v)\,(D^b\ell_j)(y',v).
\]
The legs obey $|D^a\ell_j(y,u)|\le Ce^{-\delta_2\dist(y,B(u))}$ at every site and every bond of the free lattice, $a\in\{0,1\}$ (the leg bounds \cite[(3.64)]{K1}, \cite[(2.12)]{Ba3}, at the uniform average $f=1$ in \cref{prop:a02:FL}; on the free lattice every bond is interior, so they are also \cref{lem:a01:A4}(b) and the interior case (4) of the proof of \cref{lem:a01:A5}), and $|C^{(j)}(u,v)|\le Ce^{-\delta_1|u-v|}$ (\cite[Thm.~3.3, (3.6)]{K1}; \cref{lem:a01:A23}(c)). With $|y-y'|\le\dist(y,B(u))+|u-v|+\dist(y',B(v))+\sqrt d$ and two geometric sums over $u,v\in\Z^d$, $|D^a_yD^b_{y'}G_{(j)}(y,y')|\le C_Se^{-\delta_S|y-y'|}$ at every pair of points, $\delta_S\defeq\frac12\min\{\delta_1,\delta_2\}$. In scale-$i$ units the slice $j$ has the kernel $s_j^{2-d}G_{(j)}(x/s_j,x'/s_j)$ with $s_j\defeq L^{j-i}$, and a scale-$i$ difference is $s_j^{-1}$ times a unit-$j$ one, so slice $j$ contributes at most $C_Ss_j^{-(d-2+a+b)}e^{-\delta_S|x-x'|/s_j}$; $G_1$, read in unit-$1$ units at the fixed spacing $L^{-1}$, obeys the same bound with $s\defeq L\eta$ and a constant depending on $L$ (its kernel is $L^d\ip{e_y}{G_1e_{y'}}$, bounded with the decay of \cref{lem:a01:A23}(a), and a difference at spacing $L^{-1}$ costs a factor $2L$). Finally, for $p\defeq d-2+a+b>0$ and $s$ ranging over the powers of $L$ in $[L\eta,L^{-1}]$: $e^{-\delta_S\rho/s}\le e^{-\delta_S\rho/2}e^{-\delta_S\rho/(2s)}$ as $s\le1$; the terms with $s\ge\max(\rho,\eta)$ sum to at most $(1-L^{-p})^{-1}\max(\rho,\eta)^{-p}$, and the terms with $s<\rho$ are $\rho^{-p}(\rho/s)^pe^{-\delta_S\rho/(2s)}$ with $\rho/s$ running over powers of $L$, summable to at most $C_p\rho^{-p}$. Hence $|D^a_xD^b_{x'}G_i(\eta\Z^d,0)(x,x')|\le C'(\eta+\rho)^{-p}e^{-\delta'\rho}$ with $\delta'\defeq\delta_S/2$, which is (4), and (3) gives the claim with $\delta\defeq\delta'/2$. The last clause of the statement is $(\eta+\rho)^{-p}\le(5/3)^3<5$ for $\rho\ge\frac35$.

(5) \emph{The clause (c).} By (1) every modulus is the one at $B_0=0$. The reflections of (2) and the translations of (3) are isometries of the lattice mapping bonds to bonds, so each clause of $x'\mapsto D^a_xG_i(\square,0)(x,x')$ at a pair $x',y'\in\square$ is bounded by the sum over the images of the same clause of $x'\mapsto D^a_xG_i(\eta\Z^d,0)(x,x')$ at the pair $(gx'+2\ell n,gy'+2\ell n)$, whose points are at least as far from $x$ as $x'$, $y'$ are, with the count of (3). Write $\Psi_0(x')\defeq D^a_xG_i(\eta\Z^d,0)(x,x')$ and $\rho'\defeq\min\{|x'-x|,|y'-x|\}$. \emph{Claim:} on the free lattice at zero field the sup and first-difference clauses of $\Psi_0$ at $x'$, and the H\"older quotient at a pair $x',y'$ with $|x'-y'|\le1$, are $\le Ce^{-\delta\rho'}$ whenever $\rho'\ge\frac1{20}$. Indeed, by the factorisation of (4), in unit-$j$ units $D^a_xG_{(j)}(x,\cdot)=\sum_v\ell_j(\cdot,v)\,a_j^{-1}\beta_j(v)$ is the leg of the block data $a_j^{-1}\beta_j$, $\beta_j(v)\defeq\sum_u(D^a\ell_j)(x,u)C^{(j)}(u,v)$, with $|\beta_j(v)|\le Ce^{-\delta_S\dist(x,B(v))}$ by the leg and $C^{(j)}$ bounds of (4) and one geometric sum. Split the data into the blocks with $\dist(x,B(v))\le\rho'_j/2$, $\rho'_j\defeq\rho'/s_j$ the distance in unit-$j$ units, and the rest. The leg of bounded block data obeys the sup bound of \cref{lem:a01:A4}(b), the first-difference bound of \cref{lem:a01:A5} and its H\"older bound on balls of radius $\frac1{16}$ with decay $e^{-\delta_3\dist(\cdot,S)}$ from the support $S$ of the data (the interior case, every bond of the free lattice being interior; printed: \cite[(3.63)--(3.65)]{K1}, \cite[(2.12)]{Ba3}); the far part has data of size $\le Ce^{-\delta_S\rho'_j/2}$, and the near part has its support within $\rho'_j/2$ of $x$, so at distance $\ge\rho'_j/2$ from $x'$ and $y'$; a pair at unit-$j$ distance $>\frac1{16}$ is bounded by $2\cdot16^{\alpha_0}$ times the first differences. So in unit-$j$ units every clause of slice $j$ is $\le Ce^{-\delta_6\rho'_j}$, $\delta_6\defeq\frac14\min\{\delta_S,\delta_3\}$. In scale-$i$ units the sup clause of slice $j$ scales as $s_j^{-(d-2+a)}$, the first differences as $s_j^{-(d-1+a)}$ and the H\"older quotients as $s_j^{-(d-1+a+\alpha_0)}$ (a scale-$i$ pair at distance $t\le1$ is a unit-$j$ pair at distance $t/s_j$; for $t>s_j/16$ use two first differences, $2\cdot16^{\alpha_0}s_j^{-(d-1+a)}t^{-\alpha_0}\le2\cdot16^{\alpha_0}s_j^{-(d-1+a+\alpha_0)}$); the slice $G_1$ at the fixed spacing contributes a bounded multiple of its first differences to the H\"older clause, every pair being at distance $\ge\eta$. The sum over $s$ is the geometric sum of (4) with $d-1+a+\alpha_0$ in place of $p$, giving $\le C(\eta+\rho')^{-(d-1+a+\alpha_0)}e^{-\delta'\rho'}\le Ce^{-\delta'\rho'}$ for $\rho'\ge\frac1{20}$; this is the claim, and with the image count it holds for $x'\mapsto D^a_xG_i(\square,0)(x,x')$ at pairs of $\square$ with $\rho'\ge\frac1{20}$. \emph{Localisation.} $\Psi=h'\Psi_\square$, $\Psi_\square\defeq D^a_xG_i(\square,0)(x,\cdot)$, and $D(h'\Psi_\square)=h'(\cdot+\eta e)D\Psi_\square+(\partial^\eta h')\Psi_\square$ bond by bond. Both $\Psi$ and $D\Psi$ vanish outside $S'\defeq\{\dist(\cdot,\Delta')\le\frac12+\eta\}$, whose points are at distance $\ge\frac12-2\eta\ge\frac1{10}$ from $x$ ($\eta\le\frac15$, $x$ within $\eta$ of $\Delta$, $\dist(\Delta,\Delta')\ge1$). Hence the sup and first-difference clauses of $\Psi$ are those of $\Psi_\square$ on $S'$ times $1+c_h$. For a H\"older pair $x',y'$ with both points at distance $\ge\frac1{20}$ from $x$, the Leibniz rule and the bounds $|h'(x')-h'(y')|\le dc_h|x'-y'|$, $|\partial^\eta h'(x')-\partial^\eta h'(y')|\le dc_h|x'-y'|$ give the quotient $\le C(1+c_h)^2$ from the three clauses of $\Psi_\square$; a pair with one point $y'$ within $\frac1{20}$ of $x$ has $D\Psi(y')=0$ and $|x'-y'|\ge\frac1{10}-\frac1{20}=\frac1{20}$ whenever $D\Psi(x')\ne0$, so its quotient is $\le20^{\alpha_0}$ times a first difference. This is (c).
\end{proof}
